\chapter{热力学第二定律和第三定律}
\problem{5.1}{有人想用海洋不同深度处温度的不同来制造一种机器，把海水的内能变为有用的机械功，这是否违反热力学第二定律？}
\answerarea
\problem{5.2}{欲使一密闭绝热的房间冷却，是否可以将其中的电冰箱的门打开由电冰箱运转而实现？}
\answerarea
\clearpage
\problem{5.3}{热机在工作循环中实际上与多个热源接触并交换热量。记热机从其中吸收热量的热源的最高温度为 $T_{\max}$，热机在其处放热的热源的最低温度为 $T_{\min}$，试证明该热机的效率不可能超过 $(T_{\max}-T_{\min})/T_{\max}$。}
\answerarea
\problem{5.4}{试分别利用热力学第一定律和热力学第二定律证明绝热线与等温线不能相交于两点。}
\answerarea
\clearpage
\problem{5.5}{试利用热力学第二定律证明两条绝热线不能相交。}
\answerarea
\problem{5.6}{西风吹过南北纵贯山脉，空气就会由山脉西侧的谷底越过山脉，流到山脉东侧，再向下流到与在西侧的时候同样的高度。由于高度越高，气压愈低，空气上升的时候，就会膨胀；但是并没有热量与周围的大气互换。试定性说明：
\begin{enumerate}
\item 空气流过山脉到另一侧后，温度如何变化？
\item 这样的过程是否可逆？这些空气的熵如何改变？
\end{enumerate}}
\answerarea
\clearpage
\problem{5.7}{同上题，如果空气中包含有大量的水蒸气，空气从西侧流到山顶的时候就开始凝结成雨。试定性说明：
\begin{enumerate}
\item 空气越过山脉，从西侧谷底到东侧谷底的温度和熵如何变化。
\item 这样的过程可逆吗？
\end{enumerate}}
\answerarea
\problem{5.8}{根据德拜（Debye）定律，金刚石的定体摩尔热容量随温度变化的行为是
\[
C_V=3R\frac{4\pi^4}{5}\left(\frac T\Theta\right)^3,
\]
其中 $\Theta$ 为一确定温度。现有 $1.2\,\mathrm g$ 金刚石保持体积一定，温度由 $10\,\mathrm K$ 加热到 $350\,\mathrm K$，假设形成金刚石的碳的原子量是 12，$\Theta=2230\,\mathrm K$，试以 R 为单位，确定这些金刚石的熵的改变量。}
\answerarea
\clearpage
\problem{5.9}{在一绝热材料制成的容器中，温度为 $T_1$、质量为 $m$ 的液体与质量为 $m$、温度为 $T_2$ 的同类液体在保持压强相同的条件下混合，达到平衡态。假设液体的定压比热 $c_p$ 为常量，试确定系统从初态到终态熵的变化量。}
\answerarea
\problem{5.10}{在 $24\,^\circ\mathrm C$ 时，水蒸气的饱和汽压为 $0.029824\,\mathrm{bar}$，在此条件下，求每千克水蒸气凝结为水时熵的变化。（这种情况下，每千克水蒸气凝结成水时放出热量 $2444\,\mathrm{kJ}$）}
\answerarea
\clearpage
\problem{5.11}{把 $10\,\mathrm g$、$20\,^\circ\mathrm C$ 的水在定压下变成 $-10\,^\circ\mathrm C$ 的冰，假定每克液态水的比热容量保持为 $4.2\,\mathrm{J/(g\cdot K)}$，而且冰的热容量为此值的一半，$0\,^\circ\mathrm C$ 的冰的熔解热为 $335\,\mathrm{J/g}$。试确定系统的熵的总变化量。}
\answerarea
\problem{5.12}{把 $36\,\mathrm g$、$20\,^\circ\mathrm C$ 的水，在一定的大气压强下变成 $250\,^\circ\mathrm C$ 的水蒸气。假定每克液态水的热容量保持为 $4.2\,\mathrm{J/(g\cdot K)}$，而且 $100\,^\circ\mathrm C$ 的蒸汽的潜热为 $2260\,\mathrm{J/g}$。如果蒸汽的定压摩尔热容量 $C_{p,\mathrm m}$ 与温度的关系为 $C_{p,\mathrm m}/R=a+bT+cT^2$，其中 $a=3.634$、$b=1.195\times10^{-3}\,\mathrm{K^{-1}}$、$c=0.135\times10^{-6}\,\mathrm{K^{-2}}$，试确定系统熵值的总变化量。}
\answerarea
\clearpage
\problem{5.13}{$10\,\mathrm A$ 的电流在质量为 $10\,\mathrm g$、电阻值为 $25\,\Omega$ 的电阻器上维持了一秒钟。
\begin{enumerate}
\item 如果电阻器的温度保持为 $27\,^\circ\mathrm C$，该电阻器的熵值变化了多少？电阻器与周围大气的总熵变呢？
\item 若电阻器与外界绝热，并且其定体比热容为 $0.84\,\mathrm{J/(g\cdot K)}$，开始通电时电阻器的温度为 $27\,^\circ\mathrm C$，则上述两熵变分别为多少？
\end{enumerate}}
\answerarea
\problem{5.14}{一温度为 $400\,\mathrm K$ 的热库在与另一温度为 $300\,\mathrm K$ 的热库短时间接触中传递了 $1\,\mathrm{cal}$ 的热量，两热库构成的系统的熵改变了多少？}
\answerarea
\clearpage
\problem{5.15}{初始时，均匀细杆的一端的温度为 $T_1$，另一端的温度为 $T_2$。经过一段时间后，杆上温度处处均匀，都为 $(T_1+T_2)/2$。试确定细杆在这两状态间的熵的改变量。}
\answerarea
\problem{5.16}{冬季房间热量的流失率为 $2.5\times10^4\,\mathrm{kcal/h}$，室温 $21\,^\circ\mathrm C$，外界气温 $-5\,^\circ\mathrm C$，室内外状态的熵增加率为多大？}
\answerarea
\clearpage
\problem{5.17}{试导出理想气体的摩尔熵随系统温度及体积变化的行为以及理想气体系统的动力学黏度随温度变化的行为，并给出黏度与摩尔熵的比值的表达式，画出其随温度变化的行为的曲线。}
\answerarea
\problem{5.18}{设有 $1\,\mathrm{mol}$ 理想气体从平衡态 1 变到平衡态 3，试利用图中所示的可逆过程计算其熵的变化，并与直接用理想气体熵的公式所得的结果进行比较。
\fig{5_18}}
\answerarea
\clearpage
\problem{5.19}{如图，$1\,\mathrm{mol}$ 理想气体（$\gamma=1.4$）在状态 $A$ 的参量 $V_A=20\,\mathrm L$、$T_A=300\,\mathrm K$；在状态 $C$ 的参量为 $V_C=40\,\mathrm L$、$T_C=300\,\mathrm K$。图中 $AC$ 为等温线，$AD$ 为绝热线，$AB,DC$ 为等压线，$BC$ 为等体线。试分别由下述三条路径计算 $S_C-S_A$：
\begin{enumerate}
\item $A\longrightarrow B\longrightarrow C$；
\item $A\longrightarrow C$；
\item $A\longrightarrow D\longrightarrow C$。
\end{enumerate}
\fig{5_19}}
\answerarea
\problem{5.20}{假设汽车以 $120\,\mathrm{km/h}$ 的速率在水平道路上行驶时克服空气阻力和内外各种摩擦消耗的功率为 $25\,\mathrm{kW}$，并且环境温度为 $20\,^\circ\mathrm C$，试确定其熵产生率。}
\answerarea
\clearpage
\problem{5.21}{重量为 $2100\,\mathrm{kg}$ 的汽车以 $120\,\mathrm{km/h}$ 的速率行使时突然刹车，停止时闸瓦的温度升高到 $60\,^\circ\mathrm C$，环境温度为 $20\,^\circ\mathrm C$。
\begin{enumerate}
\item 在闸瓦处机械能耗散为热时产生多少熵？
\item 在闸瓦处热量散布到空气中时产生多少附加熵？
\end{enumerate}}
\answerarea
\problem{5.22}{由 $N$ 个微观粒子组成的气体处于容积为 $V_t$ 的容器内，每个粒子都可能处于容器内的任何一处。记容器中的一部分的体积为 $V_p$（$V_p<V_t$）。
\begin{enumerate}
\item 试证明，观测到 $N$ 个粒子都处于 $V_p$ 内的概率为 $P_p=\left(\dfrac{V_p}{V_t}\right)^N$；
\item 试证明，在等温情况下，由 $N$ 个粒子组成的气体系统的体积由 $V_p$ 膨胀到 $V_t$ 后，气体系统的熵的增加量为 $\Delta S_E=-k_{\rm B}\ln P_p$，其中 $P_p=\left(\dfrac{V_p}{V_t}\right)^N$，$k_{\rm B}$ 为玻尔兹曼常量。
\end{enumerate}}
\answerarea
\clearpage
\problem{5.23}{一容器内储有由 1000 个微观粒子组成的气体，这些粒子总有机会都处于容器的左半部分，而容器的右半部分完全是空的。试计算，如果观测时间为宇宙寿命长度，即 150 亿年，则在这段时间内，可观测到上述情况的时间为多少秒？并从所得数据出发，说明自由膨胀过程是不可逆过程。}
\answerarea
\problem{5.24}{一热机在温度为 $T_1$ 的物体和温度为 $T_2$ 的热源之间循环工作，使物体的温度下降到 $T_2$ 为止。在该过程中，热机从物体吸收热量 $\Delta Q_1$，对外做功 $W'$，物体的熵的增量为 $\Delta S=S_2-S_1$。试由熵增加原理证明：
\[
W'\leq\Delta Q_1+T_2(S_2-S_1).
\]}
\answerarea
\clearpage
\problem{5.25}{一热机工作于两个温度分别为 $T_1,T_2$ 的物体之间，若工作过程中压强保持不变，两物体的热容量分别为 $C_{p,1},C_{p,2}$，试确定该热机可作的最大功。}
\answerarea
\problem{5.26}{一物体的温度和热源的温度相同，都为 $T_1$。现有一制冷机在此物体和热源之间工作，将物体的温度从 $T_1$ 降到 $T_2$，整个过程中制冷机从物体吸收热量 $Q$，熵的减少量为 $S_1-S_2$，试证明：该过程中外界所提供的功至少为 $W_{\min}=T_1(S_1-S_2)-Q$。}
\answerarea
\clearpage
\problem{5.27}{两个相同的物体的定压热容量 $C_p$ 为常量，温度都为 $T_i$，现使一制冷机在这两个物体之间工作（保持压强不变），把一个物体的温度降为 $T_f$。试证明：该过程所需的最小功为
\[
W_{\min}=C_p\left[\frac{T_i^2}{T_f}+T_f-2T_i\right].
\]}
\answerarea
\problem{5.28}{实用中也常以温度 $T$ 和熵 $S$ 作为描述均匀系统的独立状态参量，这样的以熵 $S$ 为横坐标、温度 $T$ 为纵坐标的图称为温熵图。试在温熵图上表示出卡诺循环过程，并用此图计算卡诺循环的效率。}
\answerarea
\clearpage
\problem{5.29}{试给出范德瓦尔斯气体在准静态绝热膨胀过程中的温度改变量 $\Delta T$ 随压强改变量 $\Delta p$ 的变化行为 $\Delta T/\Delta p$。}
\answerarea
\problem{5.30}{现有气体节流膨胀装置和绝热膨胀装置，试通过具体计算分析，设计一个可以使范德瓦尔斯气体高效率降温的装置。}
\answerarea
\clearpage
\problem{5.31}{试导出范德瓦尔斯模型下，$\nu\,\mathrm{mol}$ 物质的热力学熵、自由能和自由焓的表达式 $S(T,V),F(T,V),G(T,V)$，画出对应于一确定温度的 $S(V),F(V)$ 和 $G(V)$ 曲线。并与理想气体系统的相应曲线进行比较，说明其间的异同。}
\answerarea
\problem{5.32}{试证明：范德瓦尔斯气体的定体热容量只是温度的函数，与系统的体积无关。}
\answerarea
\clearpage
\problem{5.33}{实验测量表明，一弹簧的恢复力 $F_R$ 与其伸长量 $x$ 成正比，即有 $F_R=-Kx$，其中的劲度系数 $K$ 是温度的函数。试证明，如果弹簧本身的热膨胀可以忽略，则该弹簧的自由能、熵、内能可以分别表示为
\[
\begin{aligned}
F(T,x)&=F(T,0)+\frac12Kx^2,\\
S(T,x)&=S(T,0)-\frac12x^2\frac{\mathrm dK}{\mathrm dT},\\
U(T,x)&=U(T,0)+\frac12\left(K-T\frac{\mathrm dK}{\mathrm dT}\right)x^2.
\end{aligned}
\]}
\answerarea
\problem{5.34}{实验测量表明，一状态参量为压强 $p$、体积 $V$ 和温度 $T$ 的物质的状态方程可以表示为 $p=\dfrac AV T^3$，其中 $A$ 为一常量。试证明：该物质的内能可以表示为
\[
U=BT^n\ln\left(\frac V{V_0}\right)+f(T),
\]
其中 $f(T)$ 为仅依赖于温度的函数，$V_0$ 为常量，并请确定 $B$ 和 $n$。}
\answerarea
\clearpage
